\documentclass[11pt]{article}
\usepackage{parskip}
\usepackage{color}

\title{Stuffed cabbage}
\author{Nathaniel Zuk}

\begin{document}
\maketitle

\textcolor{green}{(1/26/2018)}

Based on Jasmine's mom's recipe, which we got from Lolly.

\textbf{Ingredients:}

\begin{itemize}
\item Cabbage = 1 head
\item Onions = 2 sliced + 4 Tbsp grated
\item Crushed tomatoes = 3 cups \textcolor{red}{(I used 1.5 cups, 1 small can, and this might have been enough for 10 rolls)}
\item Salt = 3 tsp \textcolor{red}{(This seems like too much)}
\item Pepper = 1/2 tsp
\item Ground beef = 1 lb
\item Cooked rice = 1 cup \textcolor{red}{(I misread the direction and used 3 Tbsp of cooked rice, but this turned out ok anyway, could have been more substantial though)}
\item Egg = 1
\item Water = 3 Tbsp
\item Honey = 3 Tbsp
\item Lemon juice = 1/4 cup
\item Raisins = 1/4 cup per batch
\end{itemize}

Boil the cabbage in lightly salted water for ~8 min.  This will make it easier to remove the cabbage leaves.

Mix together the rice, ground beef, grated onion, egg, and water. \textcolor{red}{(I added the water to the broth, not the beef mixture)}

To remove each cabbage leaf, cut the leaf closest to the stem, and carefully unwrap the leaf from the head of cabbage.  Add 1-2 spoonfuls of the filling to the center of the leaf. First, fold in on both sides along the length of the stem, then roll up perpendicular to the stem.  Repeat for the rest of the filling.  As you get closer to the center of the cabbage, you may need to reheat the cabbage to make the leaves easier to remove.  \textcolor{red}{(The original recipe called for microwaving the cabbage for 5-8 min to help remove leaves)}

Place the stuffed cabbage in a large pan (the pan must have a lid!).  Pour the crushed tomatoes, sliced onions, salt, and pepper over the cabbage rolls and spread to make the sauce evenly distributed.   Cover the pan with a lid and cook on low heat for 1.5 hrs. \textcolor{red}{(The original recipe calls for cooking this in a Dutch oven or slow cooker)}

Pour the honey and lemon juice over the cabbage and cook for another 30 min. \textcolor{blue}{I've gotten by with cooking for only 15 min after this point.} \textcolor{green}{(2/11/2018)}

\textcolor{blue}{I thought that the filling should have some salt and pepper, so I added 1 tsp of salt to the filling and the rest to the broth.  Jasmine thought the salt hadn't seeped into the filling enough because it hadn't cooked long enough.  I was concerned that there was too much salt, but the saltiness is cut by the honey and the lemon juice.  I also added about a 1/3 cup of water to the broth, which probably diluted the salt too.  More crushed tomatoes would also probably cut the intensity of the salt.}

\textcolor{blue}{I thought I would need to cook it less than specified in a pan since I thought it would cook warmer.  However, after cooking for 1 hr, the cabbage still wasn't done (soft but still a little tough).  I then cooked it for another 30 min on low heat.  The following day (after we reheated it with a microwave), Jasmine said that it turned out a lot better.  So even with on the stove, go with 1.5 hrs.  }

\textcolor{green}{(3/13/2018)}

\textcolor{red}{Heating at medium low to start, not low.}

\textcolor{blue}{Jasmine said that there wasn't enough sauce.  The rolls should be swimming in sauce.  I used 1/2 a large can of crushed tomatoes and also added 1/3 cup of water.  I think this wasn't enough water.  I may also be able to use more crushed tomatoes if I use more water.}

\textcolor{green}{(12/24/2018)}

\textcolor{blue}{Michelle and Evelyn both think that I should add raisins to the sauce in order to make it sweeter.  I will try that on the next iteration.}

\textcolor{blue}{This time I split the recipe into 2 batches, one that went in a large pot and one that went into a large pan.  I put half the crushed tomatoes and chopped onions (3 onions) in each batch.  Each batch had 2 tsp of salt and ~1/2 tsp of pepper, 2 Tbsp of honey and ~3 Tbsp lemon juice.  I kept adding water to the pot and pan until the sauce was a little short of covering the stuffed cabbage.  The pan would lose water a lot, so while cooking I needed to add more water to keep the sauce at the right level.}

\end{document}